
\begin{figure*}[htbp]
    \small
    \begin{tcolorbox}[left=3pt,right=3pt,top=3pt,bottom=3pt,title=\textbf{``Fresh \& Restart'' Correction Prompt}]
    \textbf{\#\#\# Instruction}

    You are a diligent student with a sharp analytical mind. Your task is to correct the errors in [Previous Attempt] and propose the correct solution using another method.
    
    \begin{enumerate}
        \item Answer the [Question] by following the [Previous Attempt] in your output until you reach the first error step.  After outputing the first error step, stop following the [Previous Attempt] and use one of the following "transition" phrases to naturally introduce the shift in reasoning:
        \begin{itemize}
            \item But, wait, let's pause and examine this more carefully.
            \item Wait a second, let's ensure this is right. Calculating carefully:
            \item Hmm, I want to verify this calculation. Let's go through it:
            \item Wait, this doesn't seem right. Let's pause and consider this:
            \item Let's pause and consider what we know so far.
            \item This didn't seem right. Wait, let's correct that.
            \item Wait, something seems off. Let's pause and consider what we know so far.
            \item Let's pause and consider if we've set up everything correctly.
            \item Wait a second. Is everything correct? Let me double-check.
            \item Wait, maybe there's something wrong. Let's pause and reconsider.
            \item The result looks strange, is everything correct? Let me double-check.
            \item Does this make sense? Let's rethink this.
            \item Could I have missed something? Let's pause and consider what we know so far.
        \end{itemize}

    \textbf{
    \item Then, reflect on the previous solution, and propose an alternative approach to solve the question.} 
    \item Format the final answer to meet the \textbf{output requirements of} the [Question].
    \item Don't mention [Previous Attempt] and [Ground-truth Solution] explicitly. Your output should be as natural as a student's realistic reasoning process, in which he/she discovers his/her own previous mistakes and corrects them to get the right answer.
    \item If the [Previous Attempt] is actually correct, just output the [Previous Attempt] directly and format the final answer to meet the \textbf{output requirements of} the [Question].

    \end{enumerate}
    
    [Question]
    \{question\}
    
    [Previous Attempt]
    \{pred\_solution\}
    
    [Corrected Solution]
    
\end{tcolorbox}
\caption{``Fresh \& Restart'' correction prompt.}
\label{fig:fr_prompt}
\end{figure*}